\documentclass[10pt,a4paper]{article}

\input{AEDmacros}
\usepackage{caratula} % Version modificada para usar las macros de algo1 de ~> https://github.com/bcardiff/dc-tex


\titulo{Trabajo Pr\'actico 1: Especificaci\'on y WP}
\subtitulo{Subtítulo del tp}

\fecha{\today}

\materia{Algoritmos y Estructuras de Datos}
\grupo{Grupo VOID}

\integrante{Barco, Nadia}{1594/21}{barconadia38@gmail.com}
\integrante{Sequera, Jos\'e}{637/23}{joseenriquesequera@gmail.com}
\integrante{Villarroel, Victoria}{1110/21}{vicvillarroel30@gmail.com}

% Pongan cuantos integrantes quieran

% Declaramos donde van a estar las figuras
% No es obligatorio, pero suele ser comodo
\graphicspath{{../static/}}

\begin{document}

\maketitle



\section{Soluciones}

\subsection{Predicados y auxiliares}
\vspace{0.3cm}

% Predicado ES CIUDAD
\pred{esCiudad}{ciudad: \TLista{\str \times \ent}}{ciudad_1 \geq 50000}

% Predicado SIN REPETIDOS
\pred{sinRepetidos}{ciudad: \TLista{\str \times \ent}}{\paraTodo[unalinea]{i}{\ent}{0 \leq i < \longitud{ciudad} \implicaLuego \neg (\existe[unalinea]{j}{\ent}{0 \leq {j}< \longitud{ciudad} \yLuego ciudad[i]_0 = ciudad[j]_0}}}

% Predicado MATRIZ CUADRADA
\pred{esMatrizCuadrada}{distancias: \TLista{\TLista{\ent}}}{\paraTodo[unalinea]{i}{\ent}{0 \leq {i}< \longitud{distancias} \implicaLuego {\mid}{distancias}{\mid} = {\mid}{distancias[i]{\mid}}}}

% Predicado ES MATRIZ SIMETRICA
\pred{esMatrizSimetrica}{distancias: \TLista{\TLista{\ent}}}{\paraTodo[unalinea]{i}{\ent}{0 \leq {i}< \longitud{distancias} \implicaLuego{{(\paraTodo[unalinea]{j}{\ent}{0 \leq {j}< \longitud{distancias} \implicaLuego {\mid}{distancias[i][j]}{\mid} = {\mid}{distancias[j][i]{\mid}}})}}}}

% Predicado ESTAN DENTRO DEL RANGO
\pred{estanDentroDelRango}{distancias: \TLista{\TLista{\ent}}, desde: \ent, hasta: \ent}{0 \leq desde < \longitud{distancias} \land 0 \leq hasta < \longitud{distancias}}

% Predicado ES CAMINO
\pred{esCamino}{distancias: \TLista{\TLista{\ent}}, camino: \TLista{\ent}, desde: {\ent}, hasta: {\ent}}{\longitud{camino} > {1} \yLuego camino[0] = desde \land camino[\longitud{camino} - 1] = hasta \land \paraTodo[unalinea]{i}{\ent}{0 \leq {i}< \longitud{camino} - 1 \yLuego 0 \leq camino[i] < \longitud{distancias}}}

\vspace{0.3cm}


\subsection{Procedimientos}

\begin{enumerate} \setlength\itemsep{0cm}
	\item \begin{proc}{grandesCiudades}{\In ciudades: \TLista{Ciudad}}{\TLista{Ciudad}}
		\requiere{sinRepetidos(ciudades)}
		\asegura{\paraTodo[unalinea]{i}{\ent}{0 \leq {i}<{\longitud}{res} \implicaLuego esCiudad(res[i]) \land res[i] \in ciudades}{}}
	
	\end{proc}
	\vspace{0.3cm}
	\item \begin{proc}{sumaDeHabitantes}{\In menoresDeCiudades: \TLista{Ciudad}, \In mayoresDeCiudades: \TLista{Ciudad}}{\TLista{Ciudad}}
		\requiere{\longitud{menoresDeCiudades} = \longitud{mayoresDeCiudades}}
		\requiere{sinRepetidos(menoresDeCiudades) \land sinRepetidos(mayoresDeCiudades)}
		\requiere{\paraTodo[unalinea]{i}{\ent}{0 \leq {i} < \longitud{menoresDeCiudades} \implicaLuego \\
		 \existe[unalinea]{j}{\ent}{0 \leq {j}< \longitud{mayoresDeCiudades}  \land menoresDeCiudades[j]_0 = mayoresDeCiudades[j]_0}{}}{}}
		\asegura{\longitud{res} = \longitud{menoresDeCiudades}}
		\asegura{\paraTodo[unalinea]{k}{\ent}{0 \leq k < \longitud{res} \implicaLuego \existe[unalinea]{j,h}{\ent}{0 \leq j,h < \longitud{menoresDeCiudades} \land menoresDeCiudades[j]_0 = mayoresDeCiudades[h]_0 \implica res[i] = menoresDeCiudades[j]_1 + mayoresDeCiudades[h]_1}}}
	
	\end{proc}
	\vspace{0.3cm}
	\item \begin{proc}{hayCamino}{\In distancias: \TLista{\TLista{\ent}}, \In desde: \ent, \In hasta: \ent}{\bool}
		\requiere{esMatrizCuadrada(distancias)}
		\requiere{esMatrizSimetrica(distancia)}
		\requiere{estanDentroDelRango(distancias,desde,hasta)}
		\asegura{res = \True \Iff \existe[unalinea]{camino}{\TLista{\ent}}{esCamino(distancias,camino,desde,hasta)}}
	\end{proc}
	\vspace{0.3cm}
    \item \begin{proc}{cantidadCaminosNSaltos}{\Inout conexion : \TLista{\TLista{\ent}}, \In n : \ent}{}
    \requiere{esMatrizCuadrada(conexion)}
    \requiere{esMatrizSimetrica(conexion)}
    \requiere{esMatrizConSoloUnosYCeros(conexion)}
    \requiere{n > 0}
    \requiere{conexion = C_0}
    \asegura{\\
        (\exists \ conexiones : \TLista{\TLista{\TLista{\ent}}}) \ ( \longitud{conexiones} = n \land conexiones[0] = C_0 \land \\
        \paraTodo[unalinea]{i}{\ent}{((1 \leq i < n) \yLuego \longitud{conexiones[i]} = \longitud{C_0}) \implicaLuego \\
        (esMultiplicacionDeMatrices(conexiones[i-1], conexiones[0], conexiones[i]) \yLuego \\ conexion = conexiones[n-1]))}\\
    }
    
    \pred{esMatrizConSoloUnosYCeros}{conexion : \TLista{\TLista{\ent}}}{
        \paraTodo[unalinea]{i}{\ent}{(0 \leq i < \longitud{conexion}) \implicaLuego \\
        (\paraTodo[unalinea]{j}{\ent}{(0 \leq j < \longitud{conexion}) \implicaLuego \\
        conexion[i][j] = 0 \lor conexion[i][j] = 1)}}
    }

    \pred{esMultiplicacionDeMatrices}{matrizA : \TLista{\TLista{\ent}}, matrizB : \TLista{\TLista{\ent}}, matrizC : \TLista{\TLista{\ent}}}{
        ( \longitud{matrizA} = \longitud{matrizB} \land \longitud{matrizA[0]} = \longitud{matrizA} \land \longitud{matrizB[0]} = \longitud{matrizA[0]} ) \yLuego \\
        (\paraTodo[unalinea]{i}{\ent}{(0 \leq i < \longitud{matrizA}) \yLuego \longitud{matrizA[i]} = \longitud{matrizA[0]} \implicaLuego \\
        (\paraTodo[unalinea]{j}{\ent}{(0 \leq j < \longitud{matrizB}) \yLuego \longitud{matrizB[i]} = \longitud{matrizB[0]} \implicaLuego \\
        matrizC[i][j] = \sum_{k=0}^{\longitud{matrizA}-1} matrizA[i][k] \times matrizB[k][j])) }}
    }
    \end{proc}
    \vspace{0.3cm}
    \item \begin{proc}{caminoMinimo}{\In origen : \ent, \In destino : \ent, \In distancias : \TLista{\TLista{\ent}}}{\TLista{\ent}}
    \requiere{esMatrizCuadrada(distancias)}
    \requiere{esMatrizSimetrica(distancias)}
    \requiere{estanDentroDelRango(distancias, origen, destino)}
    \asegura{\\
        (\exists \ camino : \TLista{\ent}) \ ( esCamino(distancias, camino, origen, destino) \land res = camino \land \\
        \paraTodo[unalinea]{camino'}{\TLista{\ent}}{esCamino(camino', distancias, origen, destino) \implicaLuego \\
        distanciaTotal(camino) \leq distanciaTotal(camino')}\\
    }
    
    \aux{distanciaTotal}{camino : \TLista{\ent}, distancias : \TLista{\TLista{\ent}}}{\ent}{\\
    \sum_{k=0}^{\longitud{camino}-2} distancias[camino[k]][camino[k+1]]
    }
    \end{proc}
\end{enumerate}




\section{Demostraciones de correctitud}

Especificación:

\vspace{0.3 cm}

\begin{proc}{poblacionTotal}{\In ciudades : \TLista{\ent}}{\ent}
    \requiere{\existe[unalinea]{i}{\ent}{(0 \leq i < |ciudades|) \yLuego ciudades[i].habitales >       50.000)} \land \\
    \paraTodo[unalinea]{i}{\ent}{0 \leq i < |ciudades| \implicaLuego ciudades[i].habitales \geq 0} \land \\
    \paraTodo[unalinea]{i}{\ent}{0 \leq i < j < |ciudades| \implicaLuego ciudades[i].nombre \neq ciudades[j].nombre}
    }
    \asegura{ res = \sum_{i=0}^{|ciudades|-1} ciudades[i].habitantes
    }
\end{proc}

\vspace{0.3 cm}

Implementación:

\vspace{0.3 cm}

res = 0

i = 0

while (i < ciudades.length) do

\qquad res = res + ciudades[i].habitantes

\qquad i = i +1

endwhile

\subsection{Ejercicio 1}

Para probar que el programa cumple con su especificación dad una precondición P u una postcondición Q, si siempre que el programa comienza en un estado que cumple P, el programa termina su ejecución, y en el estado final se cumple Q. Evidenciemos esto con una tripla de Hoare \{P\} S \{Q\}.

\vspace{0.3 cm}

Por el colorario de monotonia \{P\} S1, while..., S3 \{Q\}:

\begin{enumerate} \setlength\itemsep{0cm}
	\item $P \Longrightarrow wp(S1, Pc)$
	\item $Pc \Longrightarrow wp(while..., Qc)$
	\item $Qc \Longrightarrow wp(S3, Q)$
\end{enumerate}

Podemos decir que $\equiv P \Longrightarrow wp(S1; while... ; S3, Q)$ es verdadera.

\vspace{0.3 cm}

Sabemos que:

\vspace{0.3 cm}

$P \equiv (\existe[unalinea]{i}{\ent}{(0 \leq i < |ciudades|) \yLuego ciudades[i].habitales >       50.000)} \land$

\qquad $\paraTodo[unalinea]{i}{\ent}{0 \leq i < |ciudades| \implicaLuego ciudades[i].habitales \geq 0} \land$

\qquad $\paraTodo[unalinea]{i}{\ent}{0 \leq i < j < |ciudades| \implicaLuego ciudades[i].nombre \neq ciudades[j].nombre})$

\vspace{0.3 cm}

$Pc \equiv (i > 0 \land res = 0 \land P)$

\vspace{0.3 cm}

$Q \equiv \sum_{i=0}^{|ciudades|-1} ciudades[i].habitantes$


\end{document}
